% Fonte principal: abstract_en e keywords_en em tcc.conf.
\begin{abstract}
Geological mapping in Brazil is still incomplete, and traditional mineral prospecting depends on field surveys and laboratory analyses that are slow and costly. Satellite spectral data cover large areas and can support the selection of sites for investigation. This work proposes to evaluate TabPFN, a foundation model for tabular data pretrained on synthetic datasets, for mineral occurrence prediction and chemical composition prediction from satellite spectral data. To this end, a dataset will be built that links the substance recorded in the active mining processes of the National Mining Agency (ANM) and the geological and geochemical references of the Geological Survey of Brazil (SGB) to reflectance values from Sentinel-2 Level-2A images. For occurrence prediction, a single classification model will assign probabilities to the substances; for chemical composition prediction, conditional on the availability of suitable references, one regression model will be fitted for each element. TabPFN will be compared with models used in related works, such as Random Forest and XGBoost, through spatial block cross-validation across Brazil and classification and regression metrics. An ablation study will assess the contribution of each group of features. The derived datasets whose redistribution is permitted will be made available on Kaggle. The work aims to verify whether this approach can support the initial screening of areas for investigation, without replacing geological confirmation in the field.
\end{abstract}
